\chapter{Calcul littéral et identités}
\label{exercices:05}
\section*{Exercice 5.1 -- Développer et réduire}
Développer $A=(3x-2)(2x+5)-4x$. Calculer ensuite $A$ pour $x=-1$.

\section*{Exercice 5.2 -- Factoriser}
Factoriser $12x^2-8x$, $25x^2-9$ et $(x+1)(2x-3)+4(x+1)$. Vérifier une des réponses en développant.

\section*{Exercice 5.3 -- Deux programmes}
Programme A : choisir $x$, ajouter $4$, élever au carré. Programme B : élever $x$ au carré, ajouter $8x$, puis $16$. Démontrer qu'ils sont identiques. Donner leur résultat pour $x=-3$.

\section*{Exercice 5.4 -- Une propriété d'entiers}
Démontrer que $(n+2)^2-n^2$ est toujours un multiple de $4$ lorsque $n$ est entier. Calculer $101^2$ sans poser la multiplication.
